\chapter{Diseases of the Stomach}
\label{ch:stomach}
The stomach is a J-shaped muscular organ that receives food from the esophagus and initiates digestion through secretion of hydrochloric acid, pepsinogen, and intrinsic factor. It has five anatomical regions: the cardia, fundus, body, antrum, and pylorus. Diseases of the stomach include acid-related disorders, inflammatory conditions, infections, and neoplasms.

%-------------------------------------------------------------
\section{Gastric Neoplasms}
%-------------------------------------------------------------

%-------------------------------------------------------------

Gastric neoplasms represent significant gastrointestinal tumors, dominated by gastric adenocarcinoma and gastric mucosa-associated lymphoid tissue (MALT) lymphoma. Chronic inflammation driven by Helicobacter pylori infection or dietary nitrites induces a progressive cascade from chronic gastritis to intestinal metaplasia and dysplasia. Endoscopic screening, surgical gastrectomy, systemic chemotherapy, and targeted HER2/immunotherapy guide clinical management.

%-------------------------------------------------------------

%-------------------------------------------------------------

%-------------------------------------------------------------

%-------------------------------------------------------------

\label{sec:Gastric Neoplasms}
%-------------------------------------------------------------%-------------------------------------------------------------

\subsection{Gastric Adenocarcinoma}
\label{sec:Gastric Adenocarcinoma}
Gastric adenocarcinoma is the fifth most common cancer worldwide and the fourth leading cause of cancer-related death. The two histological subtypes are intestinal (associated with \textit{H.~pylori}, chronic atrophic gastritis, intestinal metaplasia, and dietary factors) and diffuse (associated with CDH1 mutations, affecting younger patients with worse prognosis). The condition results from multifactorial carcinogenesis involving \textit{H.~pylori} infection, high-salt diet, smoking, family history, and pernicious anemia. Clinical features include asymptomatic early disease detected incidentally; advanced disease presents with weight loss, epigastric pain, dysphagia, vomiting, and GI bleeding. Diagnosis is based on upper endoscopy with biopsy, with staging using TNM classification, CT imaging, and laparoscopy for peritoneal assessment. Management involves endoscopic mucosal resection or surgical gastrectomy with lymphadenectomy for early-stage disease; perioperative chemotherapy (FLOT regimen), targeted therapy (trastuzumab for HER2-positive disease), and immunotherapy (nivolumab for PD-L1-positive tumors) for advanced disease. Prognosis is poor for advanced disease but favorable for early-stage disease with appropriate treatment.

\subsection{Gastric Lymphoma}
\label{sec:Gastric Lymphoma}
Gastric lymphoma is the most common extranodal lymphoma. Mucosa-associated lymphoid tissue (MALT) lymphoma accounts for the majority and is strongly associated with \textit{H.~pylori} infection. The condition results from \textit{H.~pylori}-driven lymphoproliferation; high-grade (diffuse large B-cell) lymphoma may arise from low-grade MALT lymphoma or develop de novo. Clinical features include epigastric pain, weight loss, and GI bleeding. Diagnosis is based on endoscopy with deep biopsy and immunohistochemistry. Management involves \textit{H.~pylori} eradication therapy for low-grade MALT lymphoma; antibiotic-refractory or high-grade disease requires chemotherapy (R-CHOP) and/or radiation. Prognosis is favorable for low-grade MALT lymphoma, which may regress completely with eradication therapy alone.

\subsection{Gastrointestinal Stromal Tumor}
\label{sec:Gastrointestinal Stromal Tumor}
Gastrointestinal stromal tumor (GIST) is the most common mesenchymal neoplasm of the GI tract, arising from interstitial cells of Cajal. The stomach is the most common location (60--70\% of GISTs). The condition results from activating mutations in the KIT or PDGFRA receptor tyrosine kinase genes. Clinical features range from incidental discovery to GI bleeding, abdominal pain, or a palpable mass. Diagnosis is based on endoscopy or imaging, confirmed by biopsy showing KIT (CD117) or DOG1 positivity. Management involves surgical resection for localized disease; imatinib (a tyrosine kinase inhibitor) is used for advanced, unresectable, or metastatic disease and as adjuvant therapy for high-risk resected tumors. Prognosis has been significantly improved with targeted tyrosine kinase inhibitor therapy.

%-------------------------------------------------------------
\section{Gastritis}
%-------------------------------------------------------------

%-------------------------------------------------------------

Gastritis refers to histological inflammation of the gastric mucosa caused by infectious organisms, chemical irritants, or autoimmune destruction. Impairment of protective gastric mucosal barrier mechanisms exposes parietal and chief cells to luminal gastric acid and pepsin injury. Therapy targets underlying etiologies through H. pylori eradication regimens, acid suppression with proton pump inhibitors, and avoidance of NSAIDs.

%-------------------------------------------------------------

%-------------------------------------------------------------

%-------------------------------------------------------------

%-------------------------------------------------------------

\label{sec:Gastritis}
%-------------------------------------------------------------%-------------------------------------------------------------

\subsection{Acute Gastritis}
\label{sec:Acute Gastritis}
Acute gastritis is a self-limited inflammation of the gastric mucosa. Common causes include NSAIDs, alcohol, stress (in critically ill patients), and infections; stress-related mucosal disease (SRMD) occurs in critically ill patients, particularly those requiring mechanical ventilation or with coagulopathy. The condition results from disruption of the gastric mucosal barrier by offending agents, typically affecting the fundus and body with potential progression to erosion or ulceration. Clinical features include epigastric pain, nausea, vomiting, and hematemesis. Diagnosis is based on clinical presentation; endoscopy may be performed when indicated. Management involves removal of the offending agent and acid suppression with PPIs or H2 receptor antagonists; prophylaxis is recommended for high-risk ICU patients. Prognosis is excellent with resolution upon removal of the causative factor.

\subsection{Chronic Gastritis}
\label{sec:Chronic Gastritis}
Chronic gastritis is characterized by chronic inflammatory infiltrate in the gastric mucosa. Two main types exist: Type A (autoimmune) gastritis involves the fundus and body and is associated with pernicious anemia and increased risk of gastric carcinoid tumors; Type B (atrophic) gastritis involves the antrum and is the most common form worldwide. The condition results from autoimmune destruction of parietal cells with anti-parietal cell and anti-intrinsic factor antibodies (Type A) or \textit{H.~pylori} infection (Type B). Long-standing \textit{H.~pylori} gastritis can progress through atrophic gastritis, intestinal metaplasia, dysplasia, and gastric adenocarcinoma (Correa cascade). Clinical features include dyspepsia and potential symptoms of vitamin B12 deficiency in Type A. Diagnosis is based on endoscopy with biopsy using the Sydney system for grading. Management involves \textit{H.~pylori} eradication therapy for Type B and vitamin B12 supplementation for Type A. Prognosis depends on the subtype and the presence of precancerous changes.

\subsection{Helicobacter Pylori Infection}
\label{sec:Helicobacter Pylori Infection}
\textit{Helicobacter pylori} is a gram-negative, microaerophilic, spiral bacterium that colonizes the gastric mucosa of approximately half the world's population. It is the primary cause of chronic gastritis, peptic ulcer disease, gastric mucosa-associated lymphoid tissue (MALT) lymphoma, and gastric adenocarcinoma. The condition results from fecal-oral, oral-oral, or gastro-oral transmission, with most infections acquired in childhood; the bacterium uses urease to neutralize gastric acid, flagella to penetrate the mucus layer, and adhesins to attach to epithelial cells, with virulence factors (CagA, VacA) determining disease severity. Clinical features range from asymptomatic colonization to dyspepsia, peptic ulcer disease, and gastric malignancy. Diagnosis is based on urea breath test, stool antigen test, or endoscopic biopsy-based testing (rapid urease test, histology, culture). Management involves PPI-based triple therapy (PPI + amoxicillin + clarithromycin) or bismuth quadruple therapy for 14 days, with confirmation of eradication by urea breath test or stool antigen. Prognosis is excellent with successful eradication; untreated infection carries increased risk of gastric adenocarcinoma.

%-------------------------------------------------------------
\section{Gastroesophageal Reflux Disease}
%-------------------------------------------------------------

%-------------------------------------------------------------

%-------------------------------------------------------------

%-------------------------------------------------------------

%-------------------------------------------------------------

%-------------------------------------------------------------

\label{sec:section:Gastroesophageal Reflux Disease}
%-------------------------------------------------------------%-------------------------------------------------------------

\subsection{Gastroesophageal Reflux Disease}
\label{sec:Gastroesophageal Reflux Disease}
Gastroesophageal reflux disease (GERD) is a chronic condition characterized by symptoms of heartburn and regurgitation caused by reflux of gastric contents into the esophagus, occurring when the lower esophageal sphincter (LES) relaxes inappropriately or transiently. It affects approximately 20--30\% of adults in Western countries weekly. The condition results from transient LES relaxations, reduced basal LES pressure, impaired esophageal acid clearance, and hiatal hernia, with risk factors including obesity, pregnancy, smoking, and certain foods (fatty foods, chocolate, caffeine, alcohol). Clinical features include pyrosis (heartburn), regurgitation of acid or food, dysphagia, odynophagia, chronic cough, laryngitis, and asthma-like symptoms; complications include erosive esophagitis, peptic stricture, Barrett's esophagus (intestinal metaplasia of the esophageal mucosa, a precursor to esophageal adenocarcinoma), and esophageal adenocarcinoma. Diagnosis is primarily clinical; a therapeutic trial of proton pump inhibitors (PPIs) is both diagnostic and therapeutic. Upper endoscopy is indicated for alarm features (dysphagia, weight loss, GI bleeding, anemia) or symptoms persisting despite PPI therapy. Ambulatory pH monitoring (with impedance) is the gold standard for confirming abnormal acid exposure. Management involves lifestyle modifications (weight loss, elevation of the head of the bed, avoidance of trigger foods and late meals), antacids and alginate agents for mild symptoms, H2 receptor antagonists for nocturnal symptoms, and PPIs (omeprazole, lansoprazole, esomeprazole) as the mainstay of therapy for moderate-to-severe disease. Laparoscopic Nissen fundoplication or magnetic sphincter augmentation (LINX device) are surgical options for refractory disease or patient preference to discontinue medical therapy. Endoscopic therapies (transoral incisionless fundoplication, radiofrequency ablation) are alternative options. Prognosis is generally good with appropriate medical or surgical management.

%-------------------------------------------------------------
\section{Gastroparesis}
%-------------------------------------------------------------

%-------------------------------------------------------------

%-------------------------------------------------------------

%-------------------------------------------------------------

%-------------------------------------------------------------

%-------------------------------------------------------------

\label{sec:section:Gastroparesis}
%-------------------------------------------------------------%-------------------------------------------------------------

\subsection{Gastroparesis}
\label{sec:Gastroparesis}
Gastroparesis is delayed gastric emptying in the absence of mechanical obstruction. It affects up to 5\% of diabetic patients, and also occurs post-surgically (particularly after vagotomy or fundoplication), idiopathically, or due to medications (opioids, anticholinergics) and connective tissue disorders (scleroderma). The condition results from impaired gastric motility due to autonomic neuropathy (diabetes), vagal nerve injury, or medication effects. Clinical features include nausea, vomiting of undigested food hours after eating, early satiety, postprandial fullness, and weight loss. Diagnosis is based on gastric emptying scintigraphy (4-hour study showing >10\% retention at 4 hours). Management involves dietary modification (small, frequent, low-fat, low-fiber meals), prokinetic agents (metoclopramide or domperidone), erythromycin (motilin receptor agonist), and gastric electrical stimulation (Enterra therapy) for refractory cases; metoclopramide carries a black box warning for tardive dyskinesia. Prognosis is variable, with symptomatic improvement achievable but cure uncommon.

%-------------------------------------------------------------
\section{Peptic Ulcer Disease}
%-------------------------------------------------------------

%-------------------------------------------------------------

%-------------------------------------------------------------

%-------------------------------------------------------------

%-------------------------------------------------------------

%-------------------------------------------------------------

\label{sec:Peptic Ulcer Disease}
%-------------------------------------------------------------%-------------------------------------------------------------

\subsection{Achalasia}
\label{sec:Achalasia}
Achalasia is a primary esophageal motility disorder characterized by impaired lower esophageal sphincter (LES) relaxation and absent esophageal peristalsis. The condition results from autoimmune degeneration of inhibitory nitrergic ganglion cells in the myenteric (Auerbach) plexus of the esophagus. Clinical features include progressive dysphagia to both solids and liquids, regurgitation of undigested food, chest pain, and weight loss. Diagnosis is confirmed by esophageal manometry showing incomplete LES relaxation and aperistalsis; barium swallow demonstrates a characteristic bird's beak appearance of the distal esophagus. Management options include laparoscopic Heller myotomy with fundoplication, peroral endoscopic myotomy (POEM), pneumatic balloon dilation, and botulinum toxin injection into the LES. Prognosis is good for symptom control, though patients have a slightly increased risk of esophageal squamous cell carcinoma.

\subsection{Barrett Esophagus}
\label{sec:Barrett Esophagus}
Barrett esophagus is a metaplastic condition of the distal esophagus resulting from chronic gastroesophageal reflux disease (GERD). The condition results from chronic acid and bile exposure inducing intestinal metaplasia, where normal stratified squamous epithelium is replaced by non-ciliated columnar epithelium with goblet cells. Clinical features mirror those of GERD (heartburn, regurgitation), though some patients are asymptomatic. Diagnosis requires upper endoscopy showing salmon-pink mucosa extending proximal to the gastroesophageal junction, with histopathological confirmation of intestinal goblet cells. Management includes high-dose proton pump inhibitors (PPIs) and endoscopic surveillance to screen for dysplasia; endoscopic radiofrequency ablation or mucosal resection is indicated for high-grade dysplasia. Barrett esophagus is a major risk factor for esophageal adenocarcinoma (40-fold increased risk).

\subsection{Cholecystitis}
\label{sec:Cholecystitis}
Cholecystitis is inflammation of the gallbladder wall, classified as acute calculous (90\% of cases), acalculous (10\%), or chronic. Acute calculous cholecystitis results from persistent impaction of a gallstone in the cystic duct, causing bile stasis, mucosal ischemia, chemical irritation, and secondary bacterial superinfection (Escherichia coli, Klebsiella, Enterococcus). Acalculous cholecystitis occurs in critically ill patients (burns, trauma, major surgery, prolonged parenteral nutrition). Clinical features include severe steady right upper quadrant (RUQ) or epigastric pain radiating to the right scapula/shoulder, fever, leukocytosis, and a positive Murphy sign (inspiratory arrest on deep palpation of the RUQ). Diagnosis is supported by abdominal ultrasonography showing gallbladder wall thickening ($>4$ mm), pericholecystic fluid, and sonographic Murphy sign; HIDA scan confirms cystic duct obstruction if ultrasound is equivocal. Management includes NPO status, IV fluids, IV broad-spectrum antibiotics, and early laparoscopic cholecystectomy (within 24--72 hours). Complications include gallbladder gangrene, perforation, and emphysematous cholecystitis.

\subsection{Cholelithiasis}
\label{sec:Cholelithiasis}
Cholelithiasis is the presence of solid concretions (gallstones) within the gallbladder lumen, affecting 10--15\% of the adult population. Gallstones are classified as cholesterol stones (80\%, radiolucent, resulting from bile supersaturated with cholesterol, gallbladder hypomotility, and excess bilirubin) or pigment stones (20\%, radiopaque, divided into black stones from chronic hemolysis/cirrhosis and brown stones from biliary tract infection). Risk factors for cholesterol stones include the classic '4 Fs': Female, Fat (obesity), Fertile (multiparity/estrogen therapy), and Forty ($>40$ years). Most gallstones ($>80\%$) are asymptomatic. Symptomatic cholelithiasis presents as biliary colic: episodic right upper quadrant or epigastric dull pain precipitated by fatty meal ingestion, lasting 1--5 hours, without fever or peritoneal signs. Diagnosis is by abdominal ultrasound. Asymptomatic stones require no treatment; symptomatic biliary colic is managed with elective laparoscopic cholecystectomy. Prognosis is excellent following surgery.

\subsection{Duodenal Ulcer}
\label{sec:Duodenal Ulcer}
Duodenal ulcer is the most common type of peptic ulcer, occurring predominantly in the first part of the duodenum (duodenal bulb). \textit{Helicobacter pylori} is the primary cause in 90--95\% of cases, with NSAIDs as the second most common cause. The condition results from acid hypersecretion due to loss of negative feedback from somatostatin-producing D-cells infected by \textit{H.~pylori}. Clinical features include epigastric pain relieved by food and antacids ("hunger pain"), nocturnal pain, and bloating. Diagnosis is based on upper endoscopy. Management involves \textit{H.~pylori} eradication (triple or quadruple therapy for 14 days), PPI therapy, and smoking cessation. Prognosis is favorable; unlike gastric ulcers, duodenal ulcers rarely undergo malignant transformation.

\subsection{Esophageal Varices}
\label{sec:Esophageal Varices}
Esophageal varices are dilated submucosal veins in the lower third of the esophagus resulting from portal hypertension, most commonly caused by liver cirrhosis. The condition results from collateral venous circulation routing portal blood flow from the left gastric (coronary) vein into the systemic azygos vein due to elevated intrahepatic vascular resistance. Unruptured varices are asymptomatic; rupture causes catastrophic upper gastrointestinal hemorrhage presenting with painless hematemesis, melena, hypovolemic shock, and altered consciousness. Diagnosis is by upper gastrointestinal endoscopy showing tortuous blue submucosal vessels. Management of acute variceal hemorrhage requires hemodynamic resuscitation, IV octreotide/somatostatin (to reduce splanchnic blood flow), prophylactic IV ceftriaxone, and urgent endoscopic band ligation (EBL); transjugular intrahepatic portosystemic shunt (TIPS) is indicated for refractory bleeding. Primary prophylaxis includes non-selective beta-blockers (nadolol, propranolol) or elective endoscopic band ligation. Mortality per bleeding episode is 15--20\%.

\subsection{Gastric Ulcer}
\label{sec:Gastric Ulcer}
Gastric ulcer is a mucosal defect extending through the muscularis mucosae into the deeper layers of the gastric wall, most commonly located along the lesser curvature and antrum. \textit{Helicobacter pylori} infection is present in 60--80\% of cases and NSAID use (including aspirin and low-dose aspirin) in most others. The condition results from mucosal injury caused by \textit{H.~pylori} infection, NSAID-induced prostaglandin inhibition, and less commonly Zollinger-Ellison syndrome, Crohn's disease, or CMV infection in immunocompromised patients. Clinical features include epigastric pain aggravated by food (distinguishing it from duodenal ulcers, which improve with food), nausea, vomiting, and early satiety. Diagnosis is based on upper gastrointestinal endoscopy with biopsy to exclude malignancy (all gastric ulcers should be biopsied). Management involves \textit{H.~pylori} eradication when positive, cessation of NSAIDs, and proton pump inhibitors (PPIs) for 8 weeks with repeat endoscopy to confirm healing. Prognosis is generally good with treatment; complications include hemorrhage (most common), perforation, gastric outlet obstruction, and rarely malignant transformation.

\subsection{Peptic Ulcer Bleeding}
\label{sec:Peptic Ulcer Bleeding}
Peptic ulcer bleeding is the most common complication of peptic ulcer disease and the most frequent cause of upper gastrointestinal hemorrhage, accounting for approximately 50\% of all upper GI bleeds. The condition results from erosion of the ulcer into underlying blood vessels. Clinical features range from hematemesis (vomiting of blood, either bright red or "coffee-ground") to melena (black, tarry stools) or occult GI bleeding with iron deficiency anemia. Diagnosis is based on upper endoscopy, which identifies stigmata of recent hemorrhage including active bleeding, visible vessel, adherent clot, and flat pigmented spot. Management involves hemodynamic resuscitation (IV fluids, blood transfusion), intravenous PPI infusion (to raise gastric pH and stabilize clots), and urgent endoscopy (within 24 hours) with endoscopic therapy (injection, thermal coagulation, or hemoclips) for high-risk stigmata. Prognosis is generally favorable with intervention; re-bleeding occurs in 5--10\% of cases and may require repeat endoscopy or angiographic embolization.

\subsection{Peptic Ulcer Disease}
\label{sec:Peptic Ulcer Disease}
Peptic ulcer disease (PUD) involves mucosal defects extending through the muscularis mucosae in the stomach or duodenum. The condition results from mucosal barrier disruption caused primarily by Helicobacter pylori infection (accounting for 70--90\% of duodenal and 60--70\% of gastric ulcers) and chronic nonsteroidal anti-inflammatory drug (NSAID) use. Clinical features include epigastric pain: duodenal ulcer pain typically improves with food ingestion and recurs 2--3 hours later, whereas gastric ulcer pain is often precipitated by meals. Complications include gastrointestinal bleeding, ulcer perforation (presenting with sudden severe pain and pneumoperitoneum), and gastric outlet obstruction. Diagnosis is by upper endoscopy with mucosal biopsy for H. pylori testing. Management includes proton pump inhibitors (PPIs) and H. pylori eradication therapy (quadruple therapy with PPI, bismuth, metronidazole, tetracycline, or triple therapy). Prognosis is excellent with eradication.

\subsection{Zollinger-Ellison Syndrome}
\label{sec:Zollinger-Ellison Syndrome}
Zollinger-Ellison syndrome (ZES) is a rare disorder caused by a gastrin-secreting neuroendocrine tumor (gastrinoma), located most commonly in the duodenum or pancreas (gastrinoma triangle); approximately 25--30\% of cases occur as part of Multiple Endocrine Neoplasia Type 1 (MEN1). The condition results from autonomous hypergastrinemia, causing massive gastric acid hypersecretion and severe recurrent peptic ulceration. Clinical features include refractory or multiple peptic ulcers (often extending to the jejunum), chronic watery diarrhea, and esophagitis. Diagnosis is established by elevated fasting serum gastrin ($>1000$ pg/mL) and confirmed by secretin stimulation test (showing a rise in gastrin $>120$ pg/mL); tumor localization is achieved via somatostatin receptor scintigraphy (OctreoScan) or Ga-68 DOTATATE PET-CT. Management involves high-dose proton pump inhibitors and surgical resection of localized tumors. Prognosis depends on tumor resectability and metastatic status.

